\documentclass[12pt,a4paper]{article}


%C:/Users/97252/OneDrive - ac.sce.ac.il/Lab Doc/Measuring equipment/Scop/Figure



\usepackage{float}
\usepackage{polyglossia}
\usepackage{graphicx,wrapfig,lipsum}
\usepackage{xspace}
\usepackage{glossaries}
\usepackage{color}
\usepackage{fontspec}

\usepackage{blindtext}
\usepackage[utf8]{inputenc}
\usepackage{fancyhdr}

\usepackage{booktabs} %for Excel
\usepackage{multirow}
\usepackage{pbox}

\usepackage{hyperref}

\usepackage{siunitx}
\usepackage{comment}
\usepackage{enumitem}
\usepackage{amsmath}

\usepackage{tabularx} % for table adjustment
%
\usepackage[a4paper ,total={6in, 8in}]{geometry}
%\usepackage[showframe=true]{geometry}
% XeLaTeX!
% ============================================================ %
% HEBREW support via polyglossia %
% ============================================================ %
%\usepackage{polyglossia}
\defaultfontfeatures{Mapping=tex-text, Scale=MatchLowercase}
\setdefaultlanguage{english}
\setotherlanguage{hebrew}
\newfontfamily\hebrewfont[Script=Hebrew]{David}
% Use \begin{hebrew} block of text \end{hebrew} for paragraphs.
% Use \texthebrew{ } and \textenglish{ } for short texts.
% ============================================================ %
\newcommand{\te}[1]{\textenglish{#1}}


\newcommand{\thh}[1]{\texthebrew{#1}}
%Hyper Link Setup
\hypersetup{
	colorlinks=true,
	linkcolor=blue,
	filecolor=magenta,      
	urlcolor=blue,
}

%% DMM and link to wiki
\newcommand{\dmm}{\href{https://en.wikipedia.org/wiki/Multimeter}{\te{(DMM)}}}

\begin{document}
	\begin{hebrew}
		\title{סקופ}
		\maketitle
	\end{hebrew}


\begin{hebrew}
	\section{רקע כללי}
			מכשיר הסקופ, כשמו המלא אוסילוסקופ(שמו בעברית הוא משקף תנודות), הינו מכשיר מדידה חשוב ושימושי אשר נועד בעיקר לניתוח צורות אותות מחזוריים במישור
			הזמן ע"י הצגתם על מסך הכולל סקלת ערכים מחולקת לשנתות. מכשירי הסקופ מגיעים בשתי גרסאות: אנלוגי מבוסס מסך קטודה,
			ודיגיטאלי מבוסס מסך גביש נוזלי.
			לסקופ ישנם לפחות שני ערוצי מדידה המסומנים ב Ch1 ו Ch2 .שני ערוצים מאפשרים להציג שני אותות מחזוריים בו זמנית על אותה
			סקלה במסך. לכל ערוץ יש שני הדקים המחוברים בעזרת כבל קואקסיאלי עם מחבר הנקרא BNC.\\
			
	\textbf{חשוב מאד לציין את המגבלה העיקרית של הסקופ: קיומו של הדק אחד משותף לשני הערוצים – הדק האדמה.} \\
	
			מסיבה זו, בשימוש מקבילי של שני ערוצי הסקופ, יש להקפיד לחבר את האדמה המשותפת לשני ערוצים לאותה נקודה (אותו
			פוטנציאל) במעגל הנבדק.
			חסרון נוסף של הסקופ הוא מגבלת ערכי המתח שהוא מסוגל להציג. מכאן שיש לפעמים להשתמש במעגל הנחתת מתח לפני שמחברים
			מתחים גבוהים.
			מגבלה נוספת היא חוסר יכולת הצגת זרם בסקופ באופן ישיר. להצגת זרם דרך ענף, יש לחבר בטור לענף נגד Shunt בעל ערך קטן בכדי
			שיהיה ניתן להציג את מפל המתח עליו בעזרת הסקופ.
	
\end{hebrew}












\end{document}